% page 215 of the facsimile (image p-214 of the scan)
% block 165.1 x 259.3 mm ; left margin 36.9 mm
\pgc{15.240mm}{0.000mm}{
\l{}
\l{}
\l{}
\l{\cel{54}207}
\l{}
\l{halono. I. (\sou{astron.}) Cirklo lumoza, maxim-multa-kaze koloroza,}
\l{\cel{3}qua aparas ul-instante cirkum la disko Suna,Luna o planeta}
\l{\cel{3}kande li brilas tra atmosfero vaporoza. - II. (\sou{anat.}) Mikra}
\l{\cel{3}cirklo reda qua cirkondas la mamo-pinto.}
\l{}
\l{\sou{hal}tar. (\sou{netr}ans.) Cesar sua avanco, sua marcho, ne irar plu-}
\l{\cel{3}fore ; cesar funcionar. - DEFI.}
\l{}
\l{\sou{haltero}. (\sou{gimnastiko}). Kurta stango di qua la du extremaji finas}
\l{\cel{3}per fer-bulo, quan onu elevas ed abasas por exercar la muskuli}
\l{\cel{3}di la brakii. - FIS.}
\l{}
\l{\sou{haltiko}. (\sou{zool.}) Insekto koleoptera, plantivora, saltera, e, la}
\l{\cel{3}familio \textquotedbl{}krizomelidei\textquotedbl{}, qua devoras la planti leguma e la vi-}
\l{\cel{3}to. - L.\cel{1}\sou{altica.}}
\l{}
\l{\sou{halucinar}. (\sou{netrans}.) (\sou{patol.})\cel{2}Subisar quaza alienaco neduriva}
\l{\cel{3}dum qua onu perceptas sento, kontre ke ne existas extera ob-}
\l{\cel{3}jekto apta produktar ta sento. - DEFIRS.}
\l{}
\l{\sou{haluxo}.\cel{2}La grosa pedo-fingro.}
\l{}
\l{\sou{hamadriado.}\cel{2}(\sou{mitol}.)Nimfo di la boski, di qua la fato esas li-}
\l{\cel{3}gita ad olca di arboro, e qua naskas e mortas kun ica.}
\l{}
\l{\sou{hamako}. Telo o reto suspendita horizontale per lua du extremaji,}
\l{\cel{3}tale ke lu formacas lito portebla en qua onu su kushas o su}
\l{\cel{3}ociligas. - EFRS.}
\l{}
\l{hamstro. (\sou{zool.)}\cel{2}Mikra mamifero rodera, di Rusia e di la nordo-}
\l{\cel{3}regioni di Germania. - L.\cel{1}\sou{cricetus}. - DFRS.}
\l{}
\l{hano. (\sou{zool.}) Pultro de la klaso \textquotedbl{}galinacei\textquotedbl{}, di qua la maskulo}
\l{\cel{3}havas sur la kapo kresto reda, kun tarsi forta e garnisita ye}
\l{\cel{3}sporno hokatra, e kantas per sorto di krio klara e penetranta.}
\l{\cel{3}- DE.}
\l{}
\l{hancho. (\sou{anat}.)Singla de la du parti simetra di la homo-korpo,}
\l{\cel{3}qua salias adextere, inter la kruro e la kosti, e produktita}
\l{\cel{3}da la expansuro di la iliak-osto, ube lu formacas artiko kun}
\l{\cel{3}la femuro. - DEFIS.}
\l{}
\l{handikapar. (\sou{trans.}) I. (\sou{kavalkur-konkurso}).(En kur-konkurso}
\l{\cel{3}pri qua onu admisas kavali irge-evoza ed irge-quala) kargar}
\l{\cel{3}ula ek li per pez-masi qui esas diversa segun lia evo, fort-}
\l{\cel{3}eso(grado, por egaligar la suceso-chanci. - II. (\sou{metaf.})}
\l{\cel{3}Impozar desavantajo ad ulu. - DEFIRS.}
\l{}
\l{hangaro. - Tekto pozita sur fosti o pilastri, sub qua onu shirmas}
\l{\cel{3}veturi, ligno, e c. - DEF.}
\l{}
\l{hanso. (\sou{olim}).Korporaciono de la komercisti (exemplo : la hanso}
\l{\cel{3}Teutonika). - DEFIRS.}
\l{}
\l{haptiko. La cienco qua traktas la fenomeni pri tushado.}
}
